\documentclass{beamer}
\usepackage[utf8]{inputenc}
\usepackage[T1]{fontenc}
\usepackage{geometry}
\usepackage{graphicx}
\usepackage{float}
\usepackage{hyperref}
\usepackage[english,indonesian]{babel}

\title{Optimasi Pemrosesan Template LaTeX Menggunakan Pendekatan Berbasis AI}
\author{Academic Support Agent}
\date{\today}

\begin{document}
\maketitle

\begin{abstract}

\end{abstract}

\chapter{PENDAHULUAN}


\chapter{TINJAUAN PUSTAKA}


% Original content (commented out):
% % LaTeX-Folienvorlage erstellt von Mathias Magdowski (mathias.magdowski@ovgu.de)
% veröffentlicht unter Creative-Commons-Lizenz mit Namensnennung und Weitergabe unter gleichen Bedingungen 4.0 International (CC BY-SA 4.0, https://creativecommons.org/licenses/by-sa/4.0/deed.de)
% basierend auf der PowerPoint-Vorlage von http://bit.ly/ppt-digiPH (Stand: Anfang April 2018)

% Beamer-Klasse
\documentclass{beamer}

% Pakete laden
\input{pakete}

% Stil der Online-Tagung "Hochschule digital.innovativ" der Virtuellen Pädagogischen Hochschule
\usepackage{beamer_digiPH}

\title{\#digiPH eLecture: Bitte hier den Titel einfügen}
\author{Max Mustermann}
\institute[Lehrstuhl]{
	Name des Lehrstuhls \\
	Name des Instituts/Fachbereichs \\
	Name der Universität, Ort
}
\mode<presentation>{\keywords{Schlüsselwörter durch Komma getrennt}}
\date[09.04.2018]{Datum der Präsentation, \zB 9. April 2018}

\begin{document}

\begin{withoutheadline}
\begin{frame}
	% Schriftgröße verkleinern
	\small
	% Absatzabstand einstellen
	\setlength{\parskip}{0.75\baselineskip}
	
	Dies ist die eLectures-Präsentationsvorlage der Virtuellen PH.

	Die kommenden Folien in \LaTeX\ können und \textbf{\emph{sollen} Sie nach Wunsch adaptieren}, wir bitten Sie aber, sich alle genau anzusehen und durchzulesen! Denn:

	Auf jeder Folie finden Sie Anregungen, Hilfestellungen und Hinweise, sowohl für die inhaltlich-didaktische Vorbereitung als auch die Präsentation online. Zumindest die vorhergehende Titelfolie sollte bitte auf jeden Fall in dieser Form Verwendung finden (Wiedererkennungswert!)

	Bitte senden Sie Ihre überarbeitete PDF-Datei \textbf{vor dem gemeinsamen Testtermin} direkt an Ihre/n Co-Moderator/in. Er/sie bespricht sie mit Ihnen und lädt sie für die eLecture im Raum hoch. Der Kontakt zur Co-Moderation wird aktiv von dieser hergestellt.

	\begin{block}{Zur didaktischen Planung der Stunde:}
		Eine eLecture-Stunde geht oft sehr schnell vorbei! Gehen Sie daher bitte von max. \SI{75}{\percent} der Zeit für Ihre Präsentation aus, der Rest ist meist durch Einführung und Vorstellung, Interaktion, Rückfragen und Abschluss belegt. Je interaktiver Sie die eLecture halten, desto besser: besprechen Sie sich mit Ihrer CoModeration!
	\end{block}
\end{frame}

\begin{frame}
	\maketitle
	% \maketitle funktioniert auch im Article-Modus
	% \titlepage funktioniert nur bei Präsentationen
\end{frame}
\end{withoutheadline}

\miniframesoff
\begin{frame}[label=inhalt]{Gliederung}
	\tableofcontents
\end{frame}
\miniframeson

\section{Einleitung}

% nur für die miniframe-Navigation
\subsection*{}

\begin{frame}
	\begin{block}{Zur Didaktik: Ablauf und Lernziele zu Beginn sind sehr wichtig!\\
		\small Folgendes gilt aber auch für alle anderen Auflistungen:}
		\begin{itemize}
			\item Bulletpoints bitte maßvoll verwenden! 
			\item max. 4-5/Folie für gute Aufnahme, falls nötig in Unterabschnitte unterteilen
		\end{itemize}
	\end{block}
\end{frame}

\section{Hauptteil}

% nur für die miniframe-Navigation
\subsection*{}

\begin{frame}{Folien}
	\begin{center}
		\Large
		online immer lieber {\Huge mehr} als weniger\\
		und häufig wechseln!
	\end{center}
	\pause
	\begin{center}
		\Large
		Das bringt Abwechslung und Bewegung\\auf
        den Bildschirm und hilft,\\
        {\Huge Aufmerksamkeit} zu {\Huge halten}.
	\end{center}
\end{frame}

\begin{frame}{Animationen}
	\begin{center}
		\Large
		\ldots verursachen bei Onlinepräsentationen oft Probleme, daher gilt hier:
	\end{center}
	\begin{center}
		\Huge
weniger ist mehr.
	\end{center}
	\pause
	\begin{center}
		Und gar nicht ist noch besser
	\end{center}
\end{frame}

\begin{frame}{Farbliche Hervorhebungen im Text}
	\begin{columns}
		\begin{column}{5cm}
			\begin{center}
				\includegraphics[height=6cm]{chamaeleon_hochformat}
				
				{\tiny \textcolor{digiPH_darkorange}{Bildquelle: \url{pixabay.com}, CC-0}}
			\end{center}
		\end{column}
		\begin{column}{5cm}
			\begin{center}
				\Large
				Bitte mit Bedacht!
			\end{center}
			\pause
			\begin{center}
				Nicht mehr als ein bis zwei verschiedene Schriftfarben verwenden (Ruhe, bessere Lesbarkeit).
			\end{center}
			\begin{center}
				Tipp: Lieber durch Bilder Abwechslung schaffen!
			\end{center}
		\end{column}
	\end{columns}
\end{frame}

\begin{frame}
	\begin{center}
		\includegraphics[height=6cm]{kopfhoerer}
		
		{\tiny \textcolor{digiPH_darkorange}{Nicht die Bildquelle und die Lizenz vergessen! Bildquelle: \url{pixabay.com}, CC-0}}
	\end{center}
	\begin{center}
		Haben Sie Mut zur Vereinfachung und zu ungewohnter Positionierung und finden Sie interessante Ausschnitte.
	\end{center}
\end{frame}

\begin{frame}{Worte schaffen Bilder im Kopf. Bilder pflanzen sich fort}
	\begin{center}
		\includegraphics[height=5.5cm]{puppe}
		
		{\tiny \textcolor{digiPH_darkorange}{Bildquelle: \url{pixabay.com}, CC-0}}
	\end{center}
	\begin{center}
		Geschlechtsneutrales Formulieren ist uns wichtig: bei unseren eLectures ebenso. \textbf{In Schrift, Wort und Bild!}
	\end{center}
\end{frame}

\begin{frame}{\enquote{Stopfen} Sie Ihre Folien nicht zu voll \ldots}
	\framesubtitle{Arbeiten Sie mit Bildern: online lieber mehrere, weniger dichte Folien zeigen!}
	\begin{center}
		\begin{tikzpicture}
			\node[anchor=south west,inner sep=0] (bild) at (0,0) {\includegraphics[width=0.99\textwidth]{smartphone}};
			\begin{scope}[x=(bild.south east),y=(bild.north west)]
				% Gitternetz einzeichnen
				%\draw[help lines] (0,0) grid[xstep=0.1,ystep=0.1] (1,1);
				% Text einfügen
				\node[font=\scriptsize,align=left,color=white,text width=4cm] (text) at (0.79,0.3) {Bitte verwenden Sie ausschließlich urheberrechts-einwandfreies Bildmaterial (NoGo: Schulbücher!)\\
				Am besten also: eigenes und/oder CC-Bildmaterial, \zB von \url{www.pixabay.com} oder ähnlichen Plattformen!\\
				Bitte nie eine sichtbare Bildquelle und Lizenz vergessen.};
			\end{scope}
		\end{tikzpicture}
		
		{\tiny \textcolor{digiPH_darkorange}{Bildquelle: David Bogner, 2015, \href{https://creativecommons.org/licenses/by-nc-sa/3.0/at/}{CC BY-NC-SA}}}
	\end{center}
\end{frame}

\begin{frame}{Nicht vergessen:}
	\small Ihre eLecture wird zum Nachsehen aufgezeichnet und veröffentlicht. Deshalb muss sie inhaltlich und urheberrechtlich korrekt sein!
	
	\tiny Weiteres dazu entnehmen Sie bitte dem Infoblatt: \url{http://www.virtuelle-ph.at/selbst-electures-abhalten/infoblatt}
	
	\small Die Teilnehmenden werden über einen Disclaimer zu Beginn der eLecture darüber informiert:

	\begin{center}
		\includegraphics[height=4.5cm]{disclaimer}
		
		{\tiny \textcolor{digiPH_darkorange}{Bildquelle: Lene Kieberl, \href{https://creativecommons.org/licenses/by/3.0/at/}{CC BY}}}
	\end{center}
\end{frame}

\begin{frame}{Bei Fotos von Personen,}
	\begin{columns}
		\begin{column}{5cm}
			\small besonders aber bei Minderjährigen bitte unbedingt abklären, ob diese bzw. ihre Erziehungsberechtigten mit der Veröffentlichung einverstanden sind (Recht am eigenen Bild).
		\end{column}
		\begin{column}{5cm}
			\begin{center}
				\includegraphics[height=3cm]{details_foto_xy2}
				
				{\tiny \textcolor{digiPH_darkorange}{Bildquelle: Details Foto XY, \url{pixabay.com}, CC-0}}
			\end{center}
		\end{column}
	\end{columns}
	\begin{columns}
		\begin{column}{5cm}
			\begin{center}
				\includegraphics[width=5cm]{details_foto_xy1}
				
				{\tiny \textcolor{digiPH_darkorange}{Bildquelle: Details Foto XY, \url{pixabay.com}, CC-0}}
			\end{center}
		\end{column}
		\begin{column}{5cm}
			\begin{block}{\small Tipp:}
				\scriptsize
				Als Hilfestellung können Sie vielleicht unseren \enquote{Schummelzettel OER} mit Tipps und Quellen für verwendbare Bilder nutzen, den Sie hier zum Download vorfinden:\\
				\url{http://www.virtuelle-ph.at/schummelzettel}
			\end{block}
		\end{column}
	\end{columns}
\end{frame}

\begin{frame}{Sie möchten Software/Seiten live vorzeigen?}
	% Schriftgröße verkleinern
	\scriptsize
	% Absatzabstand einstellen
	\setlength{\parskip}{0.75\baselineskip}
	
	Beachten Sie bitte, dass Sie beim Zeigen von längeren Ausschnitten einer \textbf{Programmoberfläche} (Vorzeigen von Abläufen die nicht mehr durch das Zitatrecht abgedeckt sind) gegebenenfalls die \textbf{Erlaubnis der Herstellerfirma einholen müssen}, sofern Sie nicht Urheber/in sind bzw. keine lizenzrechtliche Erlaubnis zur Vorführung vorliegt (mehr Info: \url{http://bit.ly/vphuhr}).

	\textbf{Diese Angabe kann Ihre Co-Moderation für Sie einblenden, bitte teilen Sie sie bei Testtermin oder via Email mit.}

	\begin{block}{\scriptsize Beispiel: Live-Demo \href{GIMP}{http://www.gimp.org/},GNU General Public License 29.01.2018}
		\centering
			\includegraphics[height=4cm]{gimp_demo}
	\end{block}
\end{frame}

\begin{frame}
	\frametitle<presentation>{Gleichungen (falls nötig)}
	\begin{block}{Satz des Pythagoras:}
		\begin{equation}
			a^2 + b^2 = c^2 \label{eq:pythogoras}
		\end{equation}	\end{block}
	\begin{block}{Es folgt, dass:}
		\begin{align}
			a^2 &= c^2 - b^2 \\
			b^2 &= c^2 - a^2
		\end{align} 	
 	\end{block}
\end{frame}

\begin{frame}{Diagramme}
	\begin{figure}
		\centering
			\begin{tikzpicture}
				\begin{axis}[
					beamer,
					xlabel={Zeit, $t$ (in \si{\milli\second})},
					ylabel={Spannung, $u(t)$ (in \si{\volt})},
					xmin=0,xmax=20,
					ymin=-350,ymax=350,
					legend pos=south west,
					% Beispiel für das Einbinden einer Linie durch Angabe einer Formel
					]
					\addplot+[
						domain=0:20,
						samples=101,
					] {sin(deg(x*2*pi/20))*sqrt(2)*230};
					\addlegendentry{sinusförmig};
					\addplot+[
						domain=0:20,
						samples=101,
					] {cos(deg(x*2*pi/20))*sqrt(2)*230};
					\addlegendentry{cosinusförmig};
				\end{axis}
			\end{tikzpicture}
		\caption{Harmonischer Zeitverlauf einer Spannung mit einer Frequenz von \SI{50}{\hertz} und einem Effektivwert von \SI{230}{\volt}}
		\label{fig:u_t_sinus}
	\end{figure}
\end{frame}

\begin{frame}
	\frametitle<presentation>{Tabellen}
	\begin{table}
		\caption{Beispieltabelle}
		\label{tab:beispiel}
		\centering
			\begin{tabular}{cc}
				\toprule
				Variable & Bedeutung \\
				\midrule
				$t$ & Zeit \\
				$U$ & Spannung \\
				\bottomrule
			\end{tabular}
	\end{table}
\end{frame}

\section{Zusammenfassung}

% nur für die miniframe-Navigation
\subsection*{}

\begin{frame}
	\frametitle<presentation>{Zusammenfassung}
	\begin{block}{Ergebnisse:}
	  \begin{itemize}
	  	\item Kurzfassung der wichtigsten Ergebnisse
	  	\item Aufzeigen nicht gelöster Probleme und offener Fragen
	  	\item eventuell Ausblick auf weitere geplante Arbeiten
		\end{itemize}
	\end{block}
	\begin{block}{Fragen der Zuhörer:}
	 \begin{itemize}
	 	\item genügend Zeit für Fragen einplanen
	 	\item optional: zusätzliche Folien für wahrscheinliche Fragen
	 \end{itemize}
	\end{block}	
\end{frame}

% nur für die Miniframe-Navigation
\miniframesoff

\begin{frame}<beamer>{}
	\begin{center}
		\ldots Manege frei für Ihr Knowhow und Ihre Folien!
	\end{center}
	\vspace{3cm}
	\tiny
	PS: Am Ende Ihrer eLecture weist Ihre Co-Moderation die Teilnehmenden noch auf weitere Angebote (\zB der aktuellen Online-Tagung) und die Social-Media-Präsenzen hin. Wir bitten Sie auch um Weiterverbreitung unseres (und Ihres!) Lernangebots. Selbst schon geliked? Wir freuen uns über Likes, Kommentare und natürlich Follower!
\end{frame}

\end{document}


\end{document}